\section{Service}
\label{sec:service}

Dieser Abschnitt richtet sich an geschultes Servicepersonal.

\begin{warning}
    Arbeiten am geöffneten Gerät dürfen nur von Fachpersonen ausgeführt werden.
    Vor dem Öffnen alle Prüfleitungen und das Ladegerät entfernen.
\end{warning}

%-----------------------------------------------------------------
\subsection{Kalibrierung}
\label{sec:calibration}
Sieben Messkanäle können kalibriert werden. Jeder Kanal hat einen Nullpunktschritt (Offset)
und einen optionalen Verstärkungsschritt (Gain). Jeder Schritt wird sofort im Gerät gespeichert.

\begin{tabular}{llp{8cm}}
    \rowcolor{gray} \textbf{Kanal} & \textbf{Referenzschritt} & \textbf{Beschreibung} \\
    \texttt{V\_TERM} & \qty{0.5}{\volt}           & Spannung an den Buchsen \\
    \texttt{V\_SENS} & \qty{50}{\milli\volt}      & Sense-Spannung Vierleitermessung \\
    \texttt{V\_OUT}  & \qty{0.5}{\volt}           & Ausgangsspannung des Wandlers \\
    \texttt{V\_HV}   & \qty{10}{\volt}            & Prüfspannung Isolationsmessung \\
    \texttt{I\_OUT}  & \qty{100}{\milli\ampere}   & Ausgangsstrom \\
    \texttt{I\_ISO}  & \qty{100}{\micro\ampere}   & Ableitstrom Isolationsmessung \\
    \texttt{V\_IN}   & \qty{0.5}{\volt}           & Batteriespannung \\
\end{tabular}

\begin{steps}
    \item \disp{Settings > Calibration} öffnen, Kanal mit \enc{} wählen und \btn{OK} drücken.
    \item \disp{Set input to 0}: Null anlegen (Kurzschluss oder offen, je nach Kanal) und \btn{OK}
          drücken (\disp{OK: Zero}). \btn{ESC} bricht ab.
    \item \disp{Apply known reference}: Referenz anlegen, ihren Wert mit \enc{} einstellen und
          \btn{OK} drücken (\disp{OK: Set gain}). \btn{ESC} überspringt den Verstärkungsschritt.
    \item Erscheint \disp{Gain calibration failed!}, Referenz prüfen und wiederholen.
\end{steps}
\todo[inline]{Je Kanal: benötigte Referenzquelle, Anschluss, empfohlener Referenzwert, Betriebsart in der der Kanal aktiv ist}
\todo[inline]{Benötigte Kalibriermittel (Kalibrator, Referenzwiderstände, HV-Teiler)}
\todo[inline]{Kalibrierprotokoll / Abnahmegrenzen nach der Kalibrierung}

%-----------------------------------------------------------------
\subsection{Serviceschnittstelle (CLI)}
\label{sec:cli}
Das Gerät hat eine serielle Kommandozeile am Programmierstecker.
\begin{itemize}
    \item Anschluss: \tbc{Programmieradapter ST1-3AA}, UART, \qty{115200}{Bd}, 8N1.
    \item \texttt{help} listet alle Befehle.
\end{itemize}

\begin{longtable}{p{4cm}p{11.5cm}}
    \rowcolor{gray} \textbf{Befehl} & \textbf{Funktion} \\ \endhead
    \texttt{help}                 & Befehlsliste \\
    \texttt{pADC}, \texttt{pBMS}, \texttt{pGPIO}, \texttt{pProt} & Messwerte, \ac{BMS}-Daten, IO-Zustände, Schutzzustände ausgeben \\
    \texttt{zeroCal}, \texttt{gainCal} & Nullpunkt- / Verstärkungskalibrierung eines Kanals \\
    \texttt{setOffset}, \texttt{setGain} & Kalibrierwerte direkt schreiben \\
    \texttt{setSerial}            & Seriennummer setzen \\
    \texttt{saveEEPROM}, \texttt{readEEPROM} & Konfiguration speichern / lesen \\
    \texttt{state}                & Betriebsart wechseln \\
    \texttt{control}              & Drehgeber und Tasten über das Terminal bedienen \\
    \texttt{pRTC}, \texttt{setRTC} & Echtzeituhr lesen / stellen \\
    \texttt{flashIcons}           & Menüsymbole in den Flash-Speicher schreiben \\
    \texttt{turnOff}              & Gerät ausschalten \\
\end{longtable}
\todo[inline]{Bestätigen, welche Befehle im Release-Build enthalten sind (CLI\_ENABLED / CLI\_DEBUG), Argumente dokumentieren}

%-----------------------------------------------------------------
\subsection{Firmware-Update}
\label{sec:fwupdate}
\todo[inline]{Ablauf Firmware-Update: Werkzeuge (ST-Link, STM32CubeProgrammer), Stecker, Datei, Schritte; bleibt die Konfiguration (Kalibrierung) erhalten?}
\todo[inline]{Menüsymbole: wann ist \texttt{flashIcons} / flash\_icon\_seed.py nach einem Update nötig?}

%-----------------------------------------------------------------
\subsection{Reparatur}
\todo[inline]{Kontakt für Reparaturen, Ersatzteile, Verweis auf Schemas (doc/appendix)}
